\chapter{Introduction}
\label{chap:introduction}

\section{Background and Motivation}

Planning a trip requires more than retrieving a list of popular attractions. A
traveler may need to compare destinations, arrange transportation, select an
appropriate base, find activities that match personal interests, avoid unsafe
or unwanted experiences, and combine these decisions into a feasible itinerary.
The requirements also evolve across a conversation. A broad initial question
may be followed by references such as ``there'', ``the quieter option'', or
``the final two nights'', which cannot be interpreted correctly without
conversation context.

Online tourism information is abundant but fragmented. Official tourism
portals, destination websites, travel guides, and PDF documents describe
different aspects of a destination and use inconsistent terminology. Search
engines help users locate documents, but users must still inspect and reconcile
multiple sources. Conventional recommender systems can personalize ranked
items, but they often provide limited support for open-ended questions,
multi-turn clarification, and explanations grounded in source text.

LLMs offer a natural interface for travel planning because they can interpret
free-form requests and produce readable itineraries and comparisons. However,
an LLM's parametric knowledge may be incomplete, outdated, or impossible for a
user to verify. Fluent generation may conceal unsupported details, especially
when the request contains several independent factual requirements. These
limitations motivate the use of RAG, in which external evidence is retrieved
before generation \cite{lewis2020rag}.

Basic RAG alone does not solve the complete travel-assistance problem. A single
retrieval query may fail to cover a multi-aspect request. Exact geographic
filters can remove useful evidence when metadata is incomplete. A top-ranked
passage can be relevant to one part of the question while leaving other parts
unsupported. Personalization can also be harmful if a temporary trip constraint
is stored as a permanent preference or if a soft interest becomes a hard
retrieval filter. The proposed work therefore treats travel assistance as an
inspectable pipeline combining conversation understanding, memory, hybrid
retrieval, reranking, evidence coverage, recovery, and readiness-gated
generation.

\section{Problem Statement}

The research problem addressed by this thesis is how to design and evaluate a
personalized conversational travel assistant that generates useful responses
while remaining grounded in a bounded tourism corpus.

The problem contains four connected challenges. First, the system must convert
natural, multi-turn user language into a stable representation of intent,
operation, location, constraints, and retrieval facets. Second, it must preserve
durable user preferences without confusing them with temporary trip state.
Third, retrieval must combine semantic and lexical evidence while respecting
geographic and categorical information. Fourth, the system must recognize when
retrieved evidence is incomplete and avoid presenting unsupported details as
facts.

An evaluation based only on final-answer fluency would not reveal which of these
stages failed. The thesis therefore also addresses how to evaluate query
understanding, ranked retrieval, evidence coverage, personalization,
faithfulness, and operational efficiency within one reproducible trace.

\section{Research Questions}

This thesis is guided by the following research questions:

\begin{enumerate}[label=\textbf{RQ\arabic*:}]
    \item How can persistent user preferences and temporary conversation state
    be represented separately to support relevant personalization in multi-turn
    travel conversations?

    \item How can structured parsing, hybrid dense--lexical retrieval, metadata
    signals, and CrossEncoder reranking be combined to retrieve useful tourism
    evidence?

    \item How can requirement-level coverage checking and bounded recovery
    retrieval reduce unsupported or incomplete answers to complex travel
    queries?

    \item How should a personalized travel RAG pipeline be evaluated so that
    query-understanding, retrieval, generation, latency, and cost can be
    analyzed separately?
\end{enumerate}

\section{Research Objectives}

The principal objective is to build an end-to-end Personalized Travel Guide
Assistant for Vietnam and establish an evaluation framework appropriate to its
multi-stage behavior. The specific objectives are:

\begin{itemize}
    \item construct a cleaned, chunked, metadata-enriched tourism knowledge
    base with dense embeddings and source traceability;
    \item implement conversational rewriting and structured query parsing with
    separate intent, operation, constraints, and retrieval facets;
    \item model durable user memory separately from temporary trip state;
    \item implement bounded retrieval planning for complex requests;
    \item combine dense retrieval and BM25 using RRF and rerank selected
    candidates with a CrossEncoder;
    \item check evidence coverage, perform focused recovery retrieval, and
    assign answer readiness;
    \item generate personalized answers whose factual claims are constrained by
    selected evidence; and
    \item evaluate the system on a controlled 100-query multi-turn dataset using
    deterministic metrics, independent LLM judgments, and human confirmation.
\end{itemize}

\section{Scope}

The implemented assistant focuses on travel within Vietnam. Its corpus contains
tourism documents available during the data-construction period, and the system
does not assume live access to prices, schedules, weather, or temporary closure
information. Geographic processing supports country, region, province, city,
and place-level information where available.

The work focuses on evidence-grounded recommendation and travel information,
not on completing commercial bookings or guaranteeing real-world safety. When
the corpus does not verify time-sensitive or safety-sensitive details, the
system must state the limitation. The primary experiment uses synthetic users
and queries to control preference and conversation-state requirements; it is not
a replacement for a large real-user usability study.

\section{Contributions}

The principal contributions of this thesis are:

\begin{enumerate}[label=(\roman*)]
    \item an end-to-end tourism RAG architecture that connects knowledge-base
    construction, multi-turn query understanding, personalization, retrieval,
    evidence checking, and answer generation;

    \item a structured separation between intent and operation, explicit
    constraints and retrieval facets, positive and negative requirements, and
    persistent memory and temporary state;

    \item a hybrid retrieval pipeline combining dense retrieval, BM25, RRF,
    controlled metadata signals, filter relaxation, CrossEncoder reranking, and
    destination balancing;

    \item a bounded evidence-recovery mechanism with diagnostics that
    distinguish likely corpus gaps from selection and reranking failures; and

    \item a reproducible multi-stage evaluation protocol covering understanding,
    retrieval, final answers, latency, tokens, estimated cost, and run integrity.
\end{enumerate}

\section{Thesis Organization}

Chapter~\ref{chap:related-work} reviews travel recommendation, conversational
recommendation, RAG, hybrid retrieval, reranking, corrective retrieval, and RAG
evaluation. Chapter~\ref{chap:methodology} presents the proposed methodology.
Chapter~\ref{chap:implementation} describes the implemented system and software
components. Chapter~\ref{chap:experimental-setup} defines the dataset, metrics,
model configuration, and validity controls. Chapter~\ref{chap:results} reports
and discusses the experimental results. Chapter~\ref{chap:conclusion}
summarizes the contributions, limitations, and future work.

